\documentclass[exercice]{thedocument}%
\usepackage{texmaths-tikz}
\startdatakeys
\source{}
\cycle{}
\competencesDuSocle{}
\connaissancesRequises{}
\competencesTravaillees{}
\enddatakeys
\begin{document}%
\small La figure ci-dessous n'est pas à l'échelle.
\begin{itemize}
    \item les points M, A et S sont alignés
    \item les points M, T et H sont alignés
    \item MH = 5 cm
    \item MS = 13 cm
    \item MT = 7 cm
\end{itemize}
\begin{tikzpicture}[scale=0.2]
    \coordinate[label=below left:A](A)at(0,0);
    \coordinate[label=below right:T](T)at(16.8,0);
    \coordinate[label=above left:H](H)at(16.8,12);
    \coordinate[label=right:S](S)at(28.8,12);
    \path(16.8,7)node[anchor=south east]{M};
    \draw[fill=\currentcolor!50](A)--(T)--(H)--(S)--cycle;
    \tkzMarkRightAngle[size=1,colored](H,T,A)
    \tkzMarkRightAngle[size=1,colored](T,H,S)
\end{tikzpicture}
\begin{enumerate}
    \item Démontrer que la longueur HS est égale à 12 cm.
    \item Calculer la longueur AT.
    \item Calculer la mesure de l'angle $\widehat{\text{HMS}}$.
    \item Sachant que la longueur MT est 1,4 fois plus grande que la longueur
    HM, un élève affirme : \guillemets{L'aire du triangle MAT est 1,4 fois
    plus grande que l'aire du triangle MHS.}\vskip1pt
    Cette affirmation est-elle vraie ? On rappelle que la réponse doit être
    justifiée.
\end{enumerate}
\end{document}
